\documentclass[11pt]{article}

\usepackage[a4paper,margin=1in]{geometry}
\usepackage{amsmath,amssymb,amsthm}
\usepackage{booktabs}
\usepackage[hidelinks]{hyperref}
\usepackage{microtype}
\hypersetup{
  pdftitle={The Induced Abelian Cayley Embedding Number of P3 x P4 Is 19},
  pdfauthor={Xiaojian Zeng}
}

\newtheorem{theorem}{Theorem}
\newtheorem{lemma}{Lemma}
\newcommand{\Cay}{\operatorname{Cay}}

\title{The Induced Abelian Cayley Embedding Number of
$P_3\mathbin{\square}P_4$ Is 19}
\author{Xiaojian Zeng\\
PBC School of Finance, Tsinghua University, 100190\\
\texttt{xiaojianzeng@foxmail.com}}
\date{}

\begin{document}
\maketitle

\begin{abstract}
For a finite graph $H$, let $\eta(H)$ be the least order of a finite
abelian group whose Cayley graph contains $H$ as an induced subgraph.
Fokam Souop and Bitjoka proved the product bound
$\eta(G_1\square G_2)\leq \eta(G_1)\eta(G_2)$, which gives an explicit
order-$20$ host for $P_3\square P_4$.  We determine the exact value
\[
  \eta(P_3\square P_4)=19.
\]
The upper bound is the elementary labelling
$f(i,j)=-i-4j\pmod {19}$.  The lower bound uses regularity and edge
deletion through order $15$, followed by an exact exhaustive enumeration
of every abelian group type and every possible connection set at orders
$16$, $17$, and $18$.
\end{abstract}

\section{Context and statement}

Write $P_m$ for the path on $\{0,\ldots,m-1\}$.  If $\Gamma$ is a
finite abelian group and $S\subseteq\Gamma\setminus\{0\}$ satisfies
$S=-S$, then $\Cay(\Gamma,S)$ joins distinct $x,y$ exactly when
$y-x\in S$.  An \emph{induced abelian Cayley embedding} of $H$ is an
injection $f:V(H)\to\Gamma$ for which
\[
 uv\in E(H) \quad\Longleftrightarrow\quad f(v)-f(u)\in S
 \qquad (u\ne v).
\]
The least possible $|\Gamma|$ is denoted by $\eta(H)$.

These definitions, and the exact difference-separation criterion behind
them, were introduced in this form by Fokam Souop and Bitjoka
\cite[Definition~1 and Proposition~3]{FokamSouopBitjoka2026}.  Their
Theorem~14 and Corollary~15 imply
$\eta(P_3\square P_4)\leq\eta(P_3)\eta(P_4)=4\cdot5=20$.
In the inspected version of that paper the construction at order $20$ is
verified, while the adjacent value $\eta(P_3\square P_3)=15$ is the only
exact path-grid value in its Table~1
\cite[Section~5.1 and Table~1]{FokamSouopBitjoka2026}.  The result below
improves that particular upper bound and proves optimality.  It makes no
priority claim beyond this comparison.

Let
\[
 G=P_3\square P_4,
 \qquad V(G)=\{0,1,2\}\times\{0,1,2,3\},
\]
where two vertices are adjacent when their Manhattan distance is one.
Thus $G$ has $12$ vertices, $17$ edges, and degree multiset
$(2^4,3^6,4^2)$.

\begin{theorem}\label{thm:main}
The induced abelian Cayley embedding number of the rectangular grid is
\[
  \eta(P_3\square P_4)=19.
\]
\end{theorem}

\section{An order-19 certificate}

Take $\Gamma=\mathbb Z_{19}$ and
\[
 S=\{\pm1,\pm4\}=\{1,4,15,18\},\qquad
 f(i,j)=-i-4j\pmod {19}.
\]
In row-major order the twelve labels are
\[
 (0,15,11,7,\;18,14,10,6,\;17,13,9,5).
\]
They are distinct.  Indeed, a difference $\Delta i+4\Delta j$ has
absolute value at most $14$, so congruence to zero modulo $19$ is
equality over the integers; since $|\Delta i|\leq2$, this equality forces
$\Delta j=\Delta i=0$.

A grid edge has label difference $\pm1$ or $\pm4$.  Conversely, if two
labels have difference in $S$, the integer
$\Delta i+4\Delta j\in[-14,14]$ equals one of $\pm1,\pm4$.  With
$|\Delta i|\leq2$ and $|\Delta j|\leq3$, these four equations have only
the grid-edge solutions
\[
 (\Delta i,\Delta j)=(\pm1,0),(0,\pm1).
\]
Thus $f$ is an induced embedding.  A companion exact verifier checks all
$\binom{12}{2}=66$ unordered pairs: $17$ edges and $49$ nonedges.  This
proves $\eta(G)\leq19$ without using the search that found the labelling.

\section{Lower bounds through order 15}

Orders below $12$ are excluded by injectivity.  At order $12$ the image
would be the entire Cayley graph, but every Cayley graph is regular and
$G$ is not.

Suppose now that an order-$n$ Cayley host has degree $d=|S|$, and that
$U$ is the set of omitted host vertices.  At order $13$, deleting the
single vertex in $U$ leaves every retained degree equal to $d$ or $d-1$.
This is incompatible with the three distinct grid degrees $2,3,4$.

At order $14$, deletion of two vertices and the extreme grid degrees
force $d=4$: the maximum degree gives $d\geq4$, whereas every retained
degree is at least $d-2$, so the minimum degree gives $d\leq4$.  The host
has $14d/2=28$ edges, and deletion of two vertices removes at most
$2d=8$, leaving at least $20$ rather than $17$ edges.

At order $15$, the same degree comparison gives $4\leq d\leq5$.
The handshake lemma excludes $d=5$.  If $d=4$, the host has $30$ edges
and deletion of three vertices removes at most $3d=12$, again leaving at
least $18$ edges.  Hence no host of order at most $15$ works.

\section{The exact enumeration at orders 16--18}

It remains to exclude finitely many abelian groups.  We first reduce the
possible connection-set sizes without computation.  If $|U|=n-12$ and
$e(U)$ is the number of host edges induced by $U$, deletion gives
\begin{equation}\label{eq:deletion}
  17=\frac{nd}{2}-(n-12)d+e(U).
\end{equation}
Since $G$ has a vertex of degree $4$, necessarily $d\geq4$.  Equation
\eqref{eq:deletion} gives the complete list
\[
\begin{array}{c|c}
n&(d,e(U))\\ \hline
16&(4,1)\\
17&(4,3)\\
18&(4,5),(5,2).
\end{array}
\]
For $n=17$, $d$ is even because a nonzero inverse-closed subset of an
odd-order group is a union of two-element inverse orbits.  For $n=18$,
the equation itself rules out $d\geq6$.

The finite abelian group classification leaves precisely
\begin{align*}
n=16:\quad &\mathbb Z_{16},\ \mathbb Z_2\times\mathbb Z_8,
\ \mathbb Z_4^2,\ \mathbb Z_2^2\times\mathbb Z_4,
\ \mathbb Z_2^4;\\
n=17:\quad &\mathbb Z_{17};\\
n=18:\quad &\mathbb Z_{18},\ \mathbb Z_3\times\mathbb Z_6.
\end{align*}

For completeness, we describe the exhaustive test.  Translate a
hypothetical labelling so that the fixed corner $(0,0)$ has label zero.
Translation preserves every group difference.  For each group above the
enumerator performs the following finite steps.
\begin{enumerate}
\item Enumerate every omitted $(n-12)$-subset $U$ avoiding zero.
\item Partition the nonzero group elements into negation orbits and
      enumerate every union of orbits having one of the forced sizes $d$.
      These are exactly the legal connection sets of that size.
\item Retain only pairs $(S,U)$ with the forced $e(U)$, then only those
      for which the induced graph on $\Gamma\setminus U$ has degree
      multiset $(2^4,3^6,4^2)$ and zero has degree two.
\item Test every surviving graph for an isomorphism from $G$ that maps
      $(0,0)$ to zero.  The backtracking rejects a partial bijection only
      after an exact degree or adjacency/nonadjacency mismatch with an
      already assigned vertex.
\end{enumerate}
The normalization is lossless, the inverse-orbit construction lists
every symmetric $S$, and the final bijection search lists every injection
onto a fixed image.  Therefore a genuine embedding cannot be omitted or
pruned.

Table~\ref{tab:counts} records the terminal counts.  All group arithmetic
and graph tests use exact integers; there are no randomized or
floating-point steps.
\begin{table}[ht]
\centering
\begin{tabular}{@{}cclrrr@{}}
\toprule
$n$ & group types & connection sets & $(S,U)$ pairs & rooted tests & hits\\
\midrule
16 & 5 & $21+33+33+125+1365$ & 2,152,605 & 137,792 & 0\\
17 & 1 & $28$ & 122,304 & 3,552 & 0\\
18 & 2 & $56$ per type & 1,386,112 & 11,028 & 0\\
\bottomrule
\end{tabular}
\caption{Exact exhaustive-search counts.  A rooted test is a final
corner-fixed isomorphism check after all structural filters.}
\label{tab:counts}
\end{table}

The scripts serialize the group-by-group counts and the first witness, if
any.  As an audit of the principal false-negative risk, a separate pass
replayed every enumeration while replacing the custom final backtracker
by NetworkX~3.4.2's rooted \texttt{GraphMatcher}
\cite{HagbergSchultSwart2008}; it obtained the same terminal counts and
no isomorphisms.  This use of explicit reductions, replayable
counts, and an independent terminal check is comparable in methodology
to a separate certified-search project that independently verifies its
enumeration without discovery caches \cite{Zeng2026Figshare33318282}; that
work supplies no mathematical premise for the present theorem.

Consequently no abelian Cayley host of order $16$, $17$, or $18$ contains
$G$ as an induced subgraph.  Together with the elementary exclusions at
smaller orders and the order-$19$ certificate, this proves
Theorem~\ref{thm:main}.

\section*{Computational and AI assistance}

The exhaustive enumerators and pair verifier were written in Python and
use exact finite-group and integer arithmetic.  A separate audit pass
recomputed the certificates and used an alternative graph-isomorphism
implementation.  OpenAI Codex assisted with proof organization, code and
LaTeX drafting, and audit workflow.  The theorem depends on the explicit
arguments and reproducible finite certificates described above, not on
language-model output as evidence; responsibility for the manuscript and
its claims remains with the author.

\bibliographystyle{abbrv}
\bibliography{references}

\end{document}
